% !TeX root = ../../main.tex
% =====================================================================
%  Definición de los puntos duros
%  Responsable:
%  Última revisión:
% =====================================================================
\chapter{Definición de los puntos duros}
\label{ch:puntos_duros}

\begin{objetivos}
  \item Traducir los objetivos numéricos del bloque de dinámica vehicular en
        una lista de puntos duros defendible.
  \item Elegir herramienta con criterio y saber qué ve y qué no ve cada una.
  \item Hacer un estudio paramétrico que sirva para decidir, y no solo para
        llenar páginas del informe.
  \item Documentar y congelar la geometría de forma que dentro de dos
        temporadas siga siendo utilizable.
\end{objetivos}

\fichasesion{90 minutos}{Capítulo \ref{ch:cinematica_suspension}}%
{Herramienta de cinemática, hoja de cálculo y la plantilla de puntos duros}

Un \textbf{punto duro} es una coordenada de la que cuelga todo lo demás: el
centro de una rótula, el eje de un pivote, el centro de rueda. Este capítulo no
añade teoría al anterior; añade \emph{método}, que es lo que separa una
geometría que se puede defender de una que salió de mover puntos hasta que la
gráfica quedó bonita.

\begin{clave}
Los puntos duros son la interfaz entre departamentos. El chasis necesita saber
dónde anclar, la transmisión dónde pasan los palieres, la aerodinámica qué
altura libre hay y el equipo de fabricación qué se puede mecanizar. En cuanto
se congelan, cualquier cambio arrastra a todos: por eso se congelan tarde en el
análisis y pronto en el calendario, y por eso el capítulo termina hablando de
control de versiones y no de geometría.
\end{clave}

% ---------------------------------------------------------------------
\section{De los objetivos de dinámica a la geometría}
\label{sec:puntos_duros:01}
\index{puntos duros}\index{objetivos de diseño}

\subsection{La cadena de decisiones}

Ningún punto duro se elige por sí mismo. Cada uno sale de un objetivo que viene
de más arriba, y el trabajo consiste en recorrer esa cadena en orden
(\cref{fig:puntos:cadena}) y no saltársela.

\begin{center}
\begin{tabularx}{\textwidth}{@{}lXl@{}}
\toprule
\textbf{Objetivo} & \textbf{Qué fija} & \textbf{Viene de} \\
\midrule
$\ay$ objetivo y masa & Vía, batalla, $\hcg$ & Cap. \ref{ch:coche_como_sistema} \\
Altura de centro de balanceo & Alturas de los pivotes interiores & Cap. \ref{ch:regimen_estacionario} \\
Ganancia de caída & Longitudes y ángulos de los trapecios & Cap. \ref{ch:neumatico} \\
Bump steer nulo & Altura y longitud del tirante & Cap. \ref{ch:cinematica_suspension} \\
Recorrido de suspensión & Altura libre y posición del balancín & Cap. \ref{ch:muelles_amortiguadores} \\
Reglamento y ergonomía & Envolvente disponible & Cap. \ref{ch:chasis_reglamento} \\
\bottomrule
\end{tabularx}
\end{center}

\begin{figure}[htbp]
\centering
\begin{tikzpicture}[node distance=5mm and 12mm]
  \node[terminal] (obj) {objetivos de dinámica vehicular};
  \node[proceso, below=of obj] (env) {envolvente: reglamento,\\piloto, llanta,
        transmisión};
  \node[proceso, below=of env] (prim) {puntos primarios:\\vía, batalla,
        centro de rueda};
  \node[proceso, below=of prim] (sec) {puntos de trapecio\\y de tirante};
  \node[proceso, below=of sec] (act) {puntos de accionamiento:\\pushrod,
        balancín, amortiguador};
  \node[decision, below=of act] (ok) {¿cumple las\\curvas objetivo?};
  \node[terminal, below=of ok] (fin) {geometría congelada};
  \node[nota, right=of prim] (n1)
       {aquí ya se cierra el \SI{80}{\percent} del comportamiento};
  \node[nota, right=of ok] (n2)
       {si no, se vuelve a los puntos de trapecio, no a los primarios};

  \draw[flujo] (obj) -- (env);
  \draw[flujo] (env) -- (prim);
  \draw[flujo] (prim) -- (sec);
  \draw[flujo] (sec) -- (act);
  \draw[flujo] (act) -- (ok);
  \draw[flujo] (ok) -- node[etiq,right] {sí} (fin);
  \draw[flujo] (ok.west) -- ++(-1.6,0) |- node[etiq,left,pos=0.25] {no} (sec.west);
  \draw[aviso] (n1) -- (prim);
  \draw[aviso] (n2) -- (ok);
\end{tikzpicture}
\caption[Cadena de decisiones de los puntos duros]{El orden en que se fijan los
puntos duros. La envolvente y los puntos primarios se cierran primero porque los
condiciona el reglamento y el empaquetado; a partir de ahí se itera sobre los
puntos de trapecio hasta cumplir las curvas objetivo. Volver a tocar los puntos
primarios en esa iteración significa que la envolvente estaba mal planteada.}
\label{fig:puntos:cadena}
\end{figure}

\subsection{Restricciones que llegan de fuera}

Antes de mover un solo punto conviene tener escrita la lista de lo que no se
puede tocar, porque cada una de estas cosas ha mandado geometrías enteras a la
basura a mitad de temporada:

\begin{itemize}
  \item \textbf{Envolvente de la llanta.} Los trapecios y el pushrod tienen que
        caber dentro de la llanta en todo el recorrido \emph{y} a tope de
        volante. Es la restricción que más sorpresas da, y se comprueba con el
        modelo 3D girando la dirección, no con la vista frontal.
  \item \textbf{Reglamento.} Ancho de vía mínimo, batalla mínima, distancia al
        suelo y las reglas de sujeción y de material de los elementos de
        suspensión.
  \item \textbf{Chasis.} Los anclajes tienen que caer en nodos del bastidor. Un
        anclaje a mitad de tubo mete flexión donde no debe y se paga en rigidez
        torsional (capítulo \ref{ch:caminos_carga_nodos}).
  \item \textbf{Palieres.} El ángulo máximo de las juntas homocinéticas limita
        dónde puede estar el centro de rueda respecto al diferencial
        (capítulo \ref{ch:palieres_juntas}).
  \item \textbf{Fabricación.} Un punto duro solo existe si alguien puede
        mecanizarlo con la tolerancia necesaria y medirlo después.
\end{itemize}

% ---------------------------------------------------------------------
\section{Flujo de trabajo y herramientas}
\label{sec:puntos_duros:02}
\index{VSusp}\index{herramientas de cinemática}

\subsection{Qué ve cada herramienta}

\begin{center}
\begin{tabularx}{\textwidth}{@{}lXX@{}}
\toprule
\textbf{Herramienta} & \textbf{Para qué sirve} & \textbf{Qué no ve} \\
\midrule
Herramienta 2D (VSusp y similares) & Iterar deprisa en vista frontal: centro de
  balanceo, caída, vía & Avance, Ackermann, interferencias, bump steer real \\
Hoja de cálculo propia & Entender de verdad lo que se hace; comprobar la
  herramienta & Lo que no se haya programado \\
CAD 3D con ensamblaje cinemático & La geometría de verdad, con interferencias y
  recorridos & Nada, pero cuesta mucho más iterar \\
Multicuerpo (Adams, MotionView) & Transitorios, flexibilidades, cargas &
  Requiere datos que un primer coche no tiene \\
\bottomrule
\end{tabularx}
\end{center}

\begin{clave}
El flujo que funciona en un equipo pequeño es \textbf{2D para explorar, 3D para
confirmar}. Se itera en la herramienta rápida hasta tener una geometría
candidata, y solo entonces se monta en el CAD, donde aparecen las
interferencias y el bump steer real. Al revés ---empezar en el CAD--- se acaba
haciendo tres iteraciones en un mes en vez de treinta en una tarde.

Y sea cual sea la herramienta: \textbf{constrúyela dos veces}. Es el apartado
siguiente y es la mejor media hora que se invierte en todo el bloque.
\end{clave}

\subsection{Construirla dos veces}

Una herramienta de cinemática no avisa de un error de entrada: si se teclea una
coordenada mal, dibuja tan contenta la curva equivocada. La única defensa barata
es calcular lo mismo por dos caminos independientes y comparar.

Lo que hay que cruzar, como mínimo:

\begin{center}
\begin{tabularx}{\textwidth}{@{}lXl@{}}
\toprule
\textbf{Magnitud} & \textbf{Por qué esta} & \textbf{Tolerancia razonable} \\
\midrule
Altura de centro de balanceo & Es muy sensible a errores de altura & \SI{2}{\milli\meter} \\
Vía en reposo & Detecta errores de coordenada lateral & \SI{1}{\milli\meter} \\
Ganancia de caída & Detecta errores en las longitudes de brazo & \SI{5}{\percent} \\
Relación de instalación & Detecta errores en el accionamiento & \SI{2}{\percent} \\
\bottomrule
\end{tabularx}
\end{center}

\begin{ojo}
Cuando dos herramientas no coinciden, la reacción natural es quedarse con la
que da el número que gusta más. Es exactamente lo contrario de lo que hay que
hacer: \textbf{una discrepancia es información}, y hay que perseguirla hasta
saber de dónde sale. En nuestra experiencia, casi siempre resultó ser un error
de entrada; el resto se explicaba por la idealización 2D frente a 3D. Ninguna
diferencia se quedó sin explicar, y esa es la regla.
\end{ojo}

% ---------------------------------------------------------------------
\section{Estudios paramétricos}
\label{sec:puntos_duros:03}
\index{estudio paramétrico}\index{sensibilidad}

\subsection{Barrer un parámetro, no diez}

Un estudio paramétrico útil mueve \textbf{un} punto duro cada vez y mira qué le
pasa a las curvas objetivo. Mover varios a la vez produce gráficas
impresionantes de las que no se puede concluir nada, porque no se sabe a quién
atribuir el cambio.

El resultado que se busca no es «la geometría óptima», sino una \textbf{tabla de
sensibilidades} en el formato del capítulo \ref{ch:coche_como_sistema}:

\begin{center}
\begin{tabularx}{\textwidth}{@{}lXl@{}}
\toprule
\textbf{Si muevo\ldots} & \textbf{Le pasa esto a\ldots} & \textbf{Unidades} \\
\midrule
Altura del pivote interior inferior & Centro de balanceo & \si{\milli\meter} por \si{\milli\meter} \\
Altura del pivote interior superior & Centro de balanceo y ganancia de caída & \si{\milli\meter}, \si{\degree\per\milli\meter} \\
Longitud del trapecio superior & Ganancia de caída y variación de vía & \si{\degree\per\milli\meter} \\
Altura del anclaje del tirante & Bump steer & \si{\degree} por \si{\milli\meter} \\
Posición del balancín & Relación de instalación & --- por \si{\milli\meter} \\
\bottomrule
\end{tabularx}
\end{center}

Con esa tabla en la mano, las discusiones de diseño se acaban en dos minutos, y
en el design event se contesta a «¿por qué este punto está aquí?» con un número
en vez de con una opinión.

\begin{ojo}
Un estudio paramétrico tiene un peligro concreto: \textbf{optimizar contra un
neumático que no es el nuestro}. Si el objetivo de ganancia de caída sale de
una curva de agarre prestada, el óptimo que salga es el óptimo de otro coche.
Mientras no haya datos propios, lo honesto es elegir valores conservadores,
dejar margen de ajuste ---anclajes intercambiables, chapas de reglaje--- y
decirlo así en el informe.
\end{ojo}

% ---------------------------------------------------------------------
\section{Documentación y congelación de la geometría}
\label{sec:puntos_duros:04}
\index{congelación de diseño}\index{documentación}

\subsection{La tabla de puntos duros}

Es el documento más importante del bloque y cabe en una hoja. Tiene que llevar,
sin excepción:

\begin{enumerate}
  \item \textbf{El origen y el convenio de ejes}, dibujados, no descritos. La
        mitad de los errores entre departamentos vienen de que cada uno pone el
        origen donde le conviene.
  \item \textbf{Una fila por punto}, con nombre normalizado, las tres
        coordenadas y a qué pieza pertenece cada lado de la unión.
  \item \textbf{La postura en la que están medidas}: coche en reposo, con
        piloto, a la altura libre de diseño. Sin eso, las coordenadas no
        significan nada.
  \item \textbf{Versión y fecha}, y quién la aprobó.
\end{enumerate}

\begin{center}
\begin{tabularx}{\textwidth}{@{}llXl@{}}
\toprule
\textbf{Código} & \textbf{Punto} & \textbf{Une} & \textbf{$x$, $y$, $z$} \\
\midrule
\texttt{FLI-1} & Trapecio inf. delantero, pivote delantero & Chasis -- trapecio & \ldots \\
\texttt{FLI-2} & Trapecio inf. delantero, pivote trasero & Chasis -- trapecio & \ldots \\
\texttt{FLO}   & Rótula inferior delantera & Trapecio -- mangueta & \ldots \\
\texttt{FTO}   & Tirante exterior delantero & Tirante -- mangueta & \ldots \\
\bottomrule
\end{tabularx}
\end{center}

\subsection{Congelar de verdad}

Congelar no es dejar de tocar el archivo: es un acto explícito con fecha, con
una versión etiquetada y con un aviso a todos los departamentos. A partir de
ahí, un cambio de punto duro deja de ser una edición y pasa a ser una
\textbf{petición de cambio} que alguien tiene que aprobar sabiendo qué arrastra.

\begin{clave}
La pregunta que hay que saber contestar antes de congelar es: \textbf{si dentro
de un año alguien quiere reconstruir esta geometría, ¿tiene todo lo que
necesita?} Si la respuesta depende de que siga en el equipo la persona que la
hizo, no está documentada. El capítulo \ref{ch:documentacion_traspaso} insiste
en esto para el coche entero; en puntos duros duele especialmente, porque es
información que no se puede recuperar mirando la pieza: las coordenadas
nominales no están escritas en ningún sitio del coche.
\end{clave}

% ---------------------------------------------------------------------
%  Cierre del capítulo
% ---------------------------------------------------------------------

\begin{caso}[Dos herramientas, un resultado]
Lo que hicimos y funcionó: montamos la cinemática \textbf{en VSusp y en una hoja
de cálculo propia}, por separado, y cruzamos altura de centro de balanceo, vía y
ganancia de caída. Las diferencias salieron de pocos milímetros y pocos puntos
porcentuales. Al perseguirlas, casi todas eran errores de entrada nuestros; el
resto se explicaba por la idealización 2D frente a 3D y por redondeos de
coordenadas.

Merece la pena decir por qué esto no es perder el tiempo: una herramienta de
cinemática no avisa de una coordenada mal tecleada, y un error de geometría que
llega al taller cuesta manguetas nuevas. Media tarde de verificación cruzada es
la póliza más barata del bloque.

También está decidido dejar \textbf{margen de ajuste donde sabemos que nos
falta criterio}: el anclaje de dirección de la mangueta es intercambiable, de
modo que la longitud y la altura del brazo de dirección se pueden mover sin
fabricar nada nuevo. Es una decisión consciente de un equipo que sabe que va a
tener que corregir.

\medskip
\otrodia[Falta la tabla de puntos duros congelada, con origen, convenio, fecha
y versión, y la tabla de sensibilidades que la justifica.]
\end{caso}

\begin{ojo}
\begin{itemize}
  \item Empezar a mover puntos antes de tener escrita la envolvente.
  \item Mover varios parámetros a la vez en un estudio paramétrico.
  \item Dar coordenadas sin decir el origen, el convenio de ejes ni la postura
        del coche.
  \item Quedarse con la herramienta que da el número que gusta cuando dos no
        coinciden.
  \item Congelar la geometría en una reunión y no avisar a los demás
        departamentos.
  \item Presentar como óptimo un resultado obtenido con un neumático prestado.
  \item Guardar la geometría solo en el fichero de la herramienta. Si la
        licencia caduca o la versión cambia, se pierde.
\end{itemize}
\end{ojo}

\begin{entregable}
\begin{enumerate}
  \item La \textbf{tabla de puntos duros} completa, con origen y convenio
        dibujados, postura de medida, versión y fecha.
  \item La \textbf{tabla de sensibilidades} de los cinco o seis parámetros que
        de verdad se han barrido.
  \item El \textbf{informe de verificación cruzada} entre dos herramientas, con
        cada diferencia explicada.
  \item El \textbf{registro de cambios} de la geometría desde la congelación,
        con motivo y aprobación de cada uno.
\end{enumerate}
\end{entregable}

\begin{juez}
\begin{enumerate}
  \item ¿De qué objetivo de dinámica vehicular sale este punto duro?
  \item ¿Qué habéis barrido y qué habéis descartado?
  \item ¿Cuánto se mueve el centro de balanceo si subo este anclaje
        \SI{10}{\milli\meter}?
  \item ¿Con qué habéis verificado la geometría además de con la herramienta
        que la generó?
  \item ¿Cuándo congelasteis y qué habéis cambiado desde entonces?
  \item ¿Comprobasteis las interferencias con la dirección a tope y la
        suspensión en el extremo del recorrido?
\end{enumerate}
\end{juez}
